\section{The double and its angular coefficients}\label{sec:geometry}

We replace reflection by straight continuation on a flat surface. This
allows us to study the spatial derivatives of an eigenfunction without
imposing boundary conditions on the sides of the triangle. Throughout
Sections~\ref{sec:geometry}--\ref{sec:rigidity}, the triangle is fixed and
nondegenerate; its angles need not satisfy an irrationality assumption.

Glue two closed copies of the triangle \(Q\) along corresponding sides,
give the copies opposite orientations, and remove the three vertices.
The resulting connected oriented surface \(M\) has a flat metric. Across
an open side, a flat chart is obtained by reflecting one copy across
that side. Consequently, changes between oriented flat charts are
rotations followed by translations. The metric completion \(\overline M\)
is compact; its missing points form a set \(V\) of three cone points,
with cone angles twice the corresponding angles of \(Q\). This double
and the coefficient framework below are also used in
\cite[Section~2]{Companion}.

Let \(SM\) be the unit tangent bundle, let \(\nu\) be normalized
Liouville measure, and denote geodesic motion by \(\Psi_t\). Thus, in
a flat chart, \(\nu=d\bar A\,d\theta/(2\pi)\), where \(d\bar A\)
is probability area measure on \(M\). We first justify using this
motion as a measure-preserving flow defined modulo a null set.

A finite-time obstruction to a geodesic can occur only at a cone point.
Indeed, a unit-speed path is Cauchy as its time approaches a finite
endpoint, and its limit in \(\overline M\) permits continuation if it
is regular. There is also no unaccounted accumulation of side crossings
at a cone point. Lift the polar angle in a punctured cone to \(\R\)
and develop the lifted geodesic into the plane. Its image is a straight
line. If the path approaches the cone tip, this line ends at the
origin, so its final segment has constant polar angle and is radial.
Apart from trajectories along a seam, which
form a null set, that segment has no side crossings. On an earlier
compact regular segment crossings cannot accumulate: a flat chart at
a proposed accumulation point makes both the path and each local seam
straight. Thus every vertex hit has a finite crossing itinerary.
For each fixed finite itinerary, unfolding makes a vertex hit the
condition that a straight trajectory pass through a fixed point. For
almost every initial position this condition selects only finitely
many directions, so it has zero Liouville measure. Taking the countable
union over itineraries, vertices, and both time directions proves that
almost every geodesic exists for all real times.

Local geodesic motion is translation at fixed angle in flat charts and
preserves Liouville measure. Following a regular trajectory through
finitely many charts proves that every time map preserves \(\nu\),
with inverse \(\Psi_{-t}\) modulo null sets. Folding the two sheets to
\(Q\), and folding their tangent vectors by the corresponding local
isometries, defines a factor map to the billiard flow. It sends \(\nu\)
to \(\mu\), since the two sheets have equal area and the angular
measure is preserved by reflection. In particular, billiard
eigenfunctions lift to geodesic-flow eigenfunctions. Factor identities
are understood almost everywhere for each fixed time.

The composition operators \(u\mapsto u\circ\Psi_t\) form a strongly
continuous unitary group on \(L^2(SM,\nu)\). To see continuity at zero,
first take a continuous compactly supported function on the regular
bundle and use pointwise continuity of regular paths and dominated
convergence. Density then gives the assertion for every \(L^2\)
function. The same assertion for the billiard flow follows through
the isometric embedding given by the factor map.

In oriented flat coordinates \(x=(x_1,x_2)\), write
\(v=(\cos\theta,\sin\theta)\), and set
\[
 \pa=\tfrac12(\partial_{x_1}-i\partial_{x_2}),\qquad
 \pb=\tfrac12(\partial_{x_1}+i\partial_{x_2}).
\]
The horizontal derivatives in the forward and perpendicular directions
are
\begin{equation}\label{eq:XY}
 X=e^{i\theta}\pa+e^{-i\theta}\pb,\qquad
 Y=i e^{i\theta}\pa-i e^{-i\theta}\pb.
\end{equation}
These vector fields are globally defined on \(SM\), commute on the
regular bundle, and are formally skew adjoint for \(\nu\).
All distributional equations in this paper use compactly supported
smooth tests on \(M\) or \(SM\). No test passes through a cone point.

We use angular Fourier coefficients, as in the flat-flow framework of
Forni and Moll \cite[Sections~3.2--3.5]{FM}. For \(j\in\Z\), let
\(P_j\) be fiberwise Fourier projection, written in a chart as
\[
 P_j u=u_j(x)e^{ij\theta},\qquad
 u_j(x)=\int_0^{2\pi}u(x,\theta)e^{-ij\theta}\frac{d\theta}{2\pi}.
\]
Under an oriented change of coordinates \(z'=e^{i\phi}z+c\),
\(\theta'=\theta+\phi\), so
\[
 u'_j=e^{-ij\phi}u_j,\qquad
 (\pa u_j)'=e^{-i(j+1)\phi}\pa u_j,\qquad
 (\pb u_j)'=e^{-i(j-1)\phi}\pb u_j.
\]
Thus a coefficient is a section of a flat unitary line bundle, and its
two complex derivatives are coefficients of indices \(j+1\) and
\(j-1\), respectively. Taking Fourier coefficients in
\eqref{eq:XY} gives, also in distributions,
\begin{equation}\label{eq:modes}
 (Xu)_j=\pa u_{j-1}+\pb u_{j+1},\qquad
 (Yu)_j=i\pa u_{j-1}-i\pb u_{j+1}.
\end{equation}
The transition factors of the two terms in each sum agree.

We write \(\nabla\) for the flat spatial gradient, keeping the chart
angle fixed for a field on \(SM\). Coefficient gradients are weak
gradients in the preceding unitary trivializations. Their norms use
\(d\bar A\); bundle norms use \(\nu\). Whenever the horizontal
derivatives are square integrable, Parseval and the orthogonality of
the forward and perpendicular directions give
\begin{equation}\label{eq:gradient}
 \norm{\nabla u}_2^2
 =\norm{Xu}_2^2+\norm{Yu}_2^2
 =\sum_{j\in\Z}\norm{\nabla u_j}_2^2.
\end{equation}
For each coefficient, pointwise algebra also gives
\(\lvert\nabla u_j\rvert^2
=2(\lvert\pa u_j\rvert^2+\lvert\pb u_j\rvert^2)\).
An \(L^2\) gradient always means a finite norm on the regular surface;
it makes no assertion about traces or operator domains at a cone tip.
Inner products are complex and linear in the first entry.

Fix a real number \(\lambda\), and put \(A=X-i\lambda\), which is
again formally skew adjoint. If
\(f\circ\Psi_t=e^{i\lambda t}f\) in \(L^2\) for each \(t\), then
\(Af=0\) locally in distributions. Indeed, for a compactly supported
smooth test \(\varphi\), measure preservation gives
\[
 \ip{f}{\varphi\circ\Psi_{-t}}
 =e^{i\lambda t}\ip{f}{\varphi}.
\]
Differentiation at zero is justified on a compact regular neighborhood
of the test support and yields the claimed equation. Our analytic goal
is the following statement, whose proof occupies the next two sections.

\begin{proposition}[Transverse rigidity]\label{prop:rigidity}
Let \(M\) be the regular oriented double of a fixed nondegenerate
Euclidean triangle, and let \(\Psi_t\) be its geodesic flow on \(SM\)
with normalized Liouville measure. Suppose \(\lambda\in\R\) and
\(f\in L^\infty(SM)\) satisfy
\[
 f\circ\Psi_t=e^{i\lambda t}f
 \quad\text{in }L^2(SM)\quad\text{for every }t\in\R.
\]
Then \(f\) has square-integrable horizontal weak derivatives and
\[
 Yf=0,\qquad Xf=i\lambda f,\qquad
 \norm{\nabla f}_2^2=\lambda^2\norm{f}_2^2.
\]
\end{proposition}

The next section first controls the gradient of each angular
coefficient of a bounded local solution of \(Ah=0\).
Section~\ref{sec:rigidity} then obtains summability of these gradients
and proves the proposition.
